\scanpage{21}
\uline{Proposition 5} Soit $\underline a$ un idéal bilatère auto-adjoint dans l'algèbre de von Neumann $\mathcal A$, $\underline b$ son adhérence uniforme. Alors tout projecteur hermitien $E\in\underline b$ est dans $\underline a$. \deleted{Soit $f$ une fonction borélienne sur $\underline C$, bornée} Soit $\underline c$ une sous-algèbre auto-adjointe \uline{commutative} \unsure{uniformt} fermée de $\underline b$, identifiée à un $C_0(M)$, alors \deleted{toute fonction \unsure{continue}} $\mathcal K(M)\subset\underline a$. Si $\underline c$ est maximale, alors \deleted{toute} partie compacte $K$ de $\Omega$ est \deleted{\unsure{définie}, et pour toute} \inserted{\unsure{borélienne dans son complément}} fonction borélienne \inserted{quasi-borélienne} \inserted{bornée} \deleted{à support compact} $f$ sur $M$, \unsure{alors}, il existe une fonction \deleted{\unsure{continue} \unsure{supérieure}} \inserted{$\unsure{f}\in C_0(M)$} \deleted{\unsure{mesurable} qui est \unsure{spectralement supérieure}} \inserted{définit un opérateur qui \unsure{fait partie de} $\underline a$.}

On a $E=\lim A_i$ ($A_i\in\underline a$); d'où $E=\lim A_i^*A_i$, donc $E=\lim B_i$ ($B_i\geq0$, $B_i\in\underline a$), d'où
\[
E=\lim EB_iE,\quad\text{i.e.}\quad E=\lim C_i
\]
où $C_i\in\underline a$, $EC_iE=C_i$. \deleted{On a $C_i\leq\|C_i\|E$, d'où $\dfrac{C_i}{\|C_i\|}\leq E$, et $\|C_i\|\longrightarrow1$, donc $\dfrac{C_i}{\|C_i\|}\longrightarrow E$, d'où $E=\lim_i D_i$, $D_i\in\underline a$, $0\leq D_i\leq E$.} Mais si $\|E-\unsure{C_{i_0}|D_{i_0}}\|\leq\frac12$, on a $E\leq2\unsure{C_{i_0}|D_{i_0}}$ (décomposition spectrale) d'où $E\in\underline a$.

\margin{On se sert seulement du fait que $\underline a$ idéal dans $\underline b$.}

Pour la deuxième partie, on \deleted{peut supposer $\underline c$ maximale (\unsure{donc} $\underline c$ est une \unsure{sous-algèbre} de $\mathcal A$)} \inserted{note que} si $f\in\mathcal K^+(M)$, on a $0\leq f\leq\varphi_K\leq g$, $\varphi_K$ fonction caractéristique d'un compact $K\subset M$, $g\in\mathcal K^+(M)$. On a $U_{\varphi_K}\leq U_g\in\underline b$, d'où $U_{\varphi_K}\in\underline b$, d'où en vertu de la partie première $U_{\varphi_K}\in\underline a$, d'où $U_f\in\underline a$. Supposons $\underline c$ maximale, ($\underline c$ est alors stable par multiplication par le centre de $\mathcal A$), \deleted{soit $K$ un \unsure{compact} \unsure{quelconque} de \unsure{$M$}. $\unsure{U_{\varphi_K}\in\underline a}$}
\mspage{22}
